% ============================================================
% SECTION 1 — INTRODUCTION
% ============================================================
\section{Introduction}\label{part-1:sec:intro}

\subsection{Historical context}
\label{part-1:ssec:historical}

The three-dimensional incompressible Navier--Stokes equations
\begin{equation}\label{part-1:eq:NS-physical}
\partial_t u + (u\cdot\nabla) u = -\nabla p + \nu \Delta u + f,
\qquad \nabla\cdot u = 0,
\end{equation}
posed on either the whole space $\R^3$ or the torus $\T^3 := \R^3/\Z^3$, 
constitute the standard continuum model for the motion of a viscous, 
incompressible, homogeneous fluid. The existence of global smooth 
solutions from smooth initial data was established in two dimensions by 
Leray in 1934 \cite{Leray1934} but remains open in three dimensions. The 
Clay Mathematics Institute, in its Millennium formulation, isolated four 
statements: two positive (global regularity in $\R^3$ and in $\T^3$, 
denoted (A) and (B)) and two negative (existence of a finite-time 
singularity in $\R^3$ and in $\T^3$, denoted (C) and (D)). 

A standard criterion for breakdown is the Beale--Kato--Majda condition 
\cite{BKM1984}: the smooth solution extends past time $T$ if and only if 
$\int_0^T \norm{\curl u(t)}_{L^\infty}\dd t < \infty$. The obstruction is 
therefore the possible blow-up of a single scalar quantity; the difficulty 
is that this quantity is critical, being neither controlled by the energy 
(the natural \emph{a priori} bound) nor by any known monotone functional.

\subsection{The gap between $L^2$ and $H^s$}
\label{part-1:ssec:gap}

The energy $\tfrac12\norm{u}_{L^2}^2$ is conserved by the Euler flow 
($\nu = f = 0$) and dissipated by Navier--Stokes. In contrast, the 
$H^s$ norms, for $s>0$, satisfy the differential inequality
\begin{equation}\label{part-1:eq:Hs-inequality}
\frac{1}{2}\frac{\dd}{\dd t}\norm{u}_{H^s}^2 
\leq C_s \norm{\nabla u}_{L^\infty}\norm{u}_{H^s}^2,
\end{equation}
whose integration by Gr\"onwall requires control of the very quantity 
which one seeks to bound. This \emph{self--referential} structure is the 
core obstruction and has resisted seventy years of attempts.

Tao \cite{Tao2016} formalised the obstruction under the name of a 
\emph{supercriticality barrier}: for any system whose critical norm is 
strictly stronger than the energy, there is a heuristic mechanism which 
prevents the flow from exploiting high-frequency instabilities. Tao 
established this barrier for an \emph{averaged} variant of 
Navier--Stokes; for the genuine equations the barrier remains conjectural.

\subsection{Contribution of the present work}
\label{part-1:ssec:contribution}

In this first paper we develop the algebraic infrastructure that will be 
used in the sequels to reformulate the four Clay statements as a single 
spectral problem. Our contributions are:

\begin{enumerate}[label=(\alph*),leftmargin=*]
\item A complete derivation, from the physical equations to the 
      spectral representation, of the triadic tensor $\Tt_{\gamma\alpha\beta}$ 
      in the helical polarisation basis (\S\ref{part-1:sec:spectral}, 
      \S\ref{part-1:sec:tensor}).

\item Three exact algebraic identities, valid for every truncation 
      $K_{\max}<\infty$: the conjugate-parity relation 
      (\S\ref{part-1:ssec:parity}), the channel decoupling 
      (\S\ref{part-1:ssec:decoupling}), and the hermitian orthogonality 
      (\S\ref{part-1:ssec:hermitian}).

\item An exact factorisation of the helical symbol into an amplitude 
      and a geometric transition factor, yielding the per-triad ratio 
      $R_{kpq}$ and proving $R_{kpq}\geq 1$ with equality if and only 
      if the triadic configuration is collinear (\S\ref{part-1:sec:triad}).

\item The identification of the Lie--Poisson structure of the truncated 
      system, with the energy and helicity as Casimirs 
      (\S\ref{part-1:ssec:LiePoisson}).
\end{enumerate}

The sequels will combine these ingredients with a renormalisation flow 
argument to reduce the four Clay statements to a single question about 
the sign of a scalar coherence functional.

\subsection{Organisation of the paper}
\label{part-1:ssec:org}

Section \ref{part-1:sec:spectral} derives the spectral formulation of 
Navier--Stokes from the physical equations. Section \ref{part-1:sec:tensor} 
introduces the triadic tensor and its eight helicoidal channels. 
Section \ref{part-1:sec:identities} proves the three algebraic identities. 
Section \ref{part-1:sec:triad} establishes the per-triad ratio and its lower 
bound. Section \ref{part-1:sec:structure} identifies the Lie--Poisson 
structure. Appendix \ref{part-1:app:proofs} collects the technical proofs.

% ============================================================
% SECTION 2 — SPECTRAL FORMULATION
% ============================================================
\section{Spectral formulation of Navier--Stokes}
\label{part-1:sec:spectral}

\subsection{The equations on the torus}
\label{part-1:ssec:torus}

We work on the three-dimensional torus $\T^3 = \R^3/\Z^3$. The 
velocity field $u:\T^3\times[0,T)\to\R^3$ and the pressure 
$p:\T^3\times[0,T)\to\R$ satisfy
\begin{equation}\label{part-1:eq:NS-torus}
\begin{cases}
\partial_t u + (u\cdot\nabla)u = -\nabla p + \nu\Delta u + f, \\[2pt]
\nabla\cdot u = 0, \\[2pt]
u(\cdot,0) = u_0,
\end{cases}
\end{equation}
with $\nu \geq 0$ and $f$ an admissible external force. We assume
\begin{equation}\label{part-1:eq:function-spaces}
u_0 \in H^s_{\mathrm{df}}(\T^3), \quad s > 5/2,
\qquad
f \in L^\infty([0,T];H^s_{\mathrm{df}}(\T^3)),
\end{equation}
where $H^s_{\mathrm{df}}$ denotes the closure of the divergence-free 
$H^s$ fields. For $s>5/2$, the velocity field is continuously embedded 
in $C^1(\T^3)$.

\begin{notation}\label{part-1:not:torus}
Throughout, $\T^3$ denotes the flat torus of side length $1$. Its 
Pontryagin dual is $\Z^3$, and we denote by $\ZZ^3_\ast := \Z^3\setminus\{0\}$ 
the set of non-zero wavevectors. All integrals are with respect to the 
normalised Lebesgue measure on $\T^3$.
\end{notation}

\subsection{The Leray projection}
\label{part-1:ssec:leray}

The Leray projection $\Leray$ onto divergence-free fields is the 
$L^2$-orthogonal projection
\begin{equation}\label{part-1:eq:Leray}
\Leray := I - \nabla \Delta^{-1} \nabla\cdot
\end{equation}
acting on $L^2(\T^3;\R^3)$, where $\Delta^{-1}$ denotes the inverse 
Laplacian on the mean-zero subspace. Its principal properties are:
\begin{enumerate}[label=(P\arabic*),leftmargin=*]
\item $\Leray^2 = \Leray$ (idempotent);
\item $\Leray^* = \Leray$ (self-adjoint on $L^2$);
\item $\ker \Leray = \{\nabla \phi : \phi \in H^1(\T^3)\}$;
\item $\operatorname{Im}\Leray = \{v\in L^2 : \nabla\cdot v = 0\}$.
\end{enumerate}

Applying $\Leray$ to \eqref{part-1:eq:NS-torus} and using 
$\Leray \nabla p = 0$, we obtain the projected evolution equation
\begin{equation}\label{part-1:eq:NS-projected}
\partial_t u = -\Leray\bigl[(u\cdot\nabla)u\bigr] + \nu\Delta u + \Leray f.
\end{equation}

\begin{lemma}[Divergence-free identity]\label{part-1:lem:div-free-identity}
For any divergence-free $u$,
\begin{equation}\label{part-1:eq:div-identity}
(u\cdot\nabla)u = \nabla\cdot(u\otimes u).
\end{equation}
Consequently
$\Leray\bigl[(u\cdot\nabla)u\bigr] = \Leray\bigl[\nabla\cdot(u\otimes u)\bigr]$.
\end{lemma}

\begin{proof}
Since $\nabla\cdot u = 0$, we have
\begin{equation*}
(\nabla\cdot(u\otimes u))_i 
= \partial_j(u_i u_j) 
= u_j \partial_j u_i + u_i \partial_j u_j 
= (u\cdot\nabla)u_i.\qedhere
\end{equation*}
\end{proof}

\subsection{Fourier transform}
\label{part-1:ssec:Fourier}

We use the Fourier convention
\begin{equation}\label{part-1:eq:Fourier}
\hat u(k,t) := \int_{\T^3} u(x,t) \e^{-2\pi \ii k\cdot x}\dd x,
\qquad k\in\Z^3,
\end{equation}
with inversion
\begin{equation}\label{part-1:eq:Fourier-inverse}
u(x,t) = \sum_{k\in\Z^3} \hat u(k,t)\, \e^{2\pi \ii k\cdot x}.
\end{equation}

The velocity field being real, the Fourier coefficients satisfy the 
reality condition
\begin{equation}\label{part-1:eq:reality}
\hat u(-k,t) = \overline{\hat u(k,t)}, \qquad k\in\Z^3.
\end{equation}

The differential operators admit the spectral symbols
\begin{equation}\label{part-1:eq:symbols}
\widehat{\nabla\cdot u}(k) = 2\pi\ii\, k\cdot\hat u(k),
\qquad
\widehat{\Delta u}(k) = -(2\pi|k|)^2 \hat u(k),
\end{equation}
and the Leray projection acts diagonally,
\begin{equation}\label{part-1:eq:Leray-symbol}
\widehat{\Leray v}(k) = P(k) \hat v(k),
\qquad
P(k) := I - \frac{k\otimes k}{\abs{k}^2}.
\end{equation}

In Fourier variables, equation \eqref{part-1:eq:NS-projected} becomes
\begin{equation}\label{part-1:eq:NS-Fourier}
\partial_t \hat u(k) 
= -2\pi\ii P(k)\!\!\sum_{p+q=k}\!\! \bigl[\hat u(p)\cdot q\bigr]\hat u(q)
- \nu(2\pi|k|)^2 \hat u(k) + P(k) \hat f(k),
\end{equation}
where the sum runs over all pairs $(p,q)\in\Z^3\times\Z^3$ with 
$p+q=k$.

\subsection{Real polarisation basis}
\label{part-1:ssec:real-pol}

For each $k\in\ZZ^3_\ast$, we construct an orthonormal basis of 
$k^\perp\subset\R^3$ as follows. Fix the reference vector 
$z := (0,0,1)$ and set $\hat k := k/|k|$. Define
\begin{equation}\label{part-1:eq:e1e2}
e_1(k) := \frac{z - (z\cdot\hat k)\hat k}{\abs{z - (z\cdot\hat k)\hat k}},
\qquad
e_2(k) := \hat k \times e_1(k).
\end{equation}
The pair $\{e_1(k),e_2(k)\}$ is orthonormal and satisfies
\begin{equation}\label{part-1:eq:polarisation-properties}
e_i(k)\cdot k = 0, \qquad
e_i(k)\cdot e_j(k) = \delta_{ij}, \qquad
e_i(-k) = e_i(k).
\end{equation}

The three-dimensional basis $\{e_1(k), e_2(k), \hat k\}$ is 
right-handed for every $k\in\ZZ^3_\ast$.

\subsection{Helicoidal basis}
\label{part-1:ssec:helical}

The \emph{helicoidal polarisation vectors} are defined by
\begin{equation}\label{part-1:eq:helicoidal}
e_+(k) := \frac{e_1(k) + \ii\, e_2(k)}{\sqrt 2},
\qquad
e_-(k) := \frac{e_1(k) - \ii\, e_2(k)}{\sqrt 2}.
\end{equation}

\begin{proposition}[Properties of the helicoidal basis]
\label{part-1:prop:helical-props}
For every $k\in\ZZ^3_\ast$, the vectors $e_\pm(k)$ satisfy:
\begin{enumerate}[label=(\roman*),leftmargin=*]
\item Hermitian orthonormality:
      $\overline{e_\sigma(k)}\cdot e_{\sigma'}(k) = \delta_{\sigma\sigma'}$;
\item Transversality: $e_\sigma(k)\cdot k = 0$;
\item Diagonalisation of the curl:
      $\ii\, k\times e_\sigma(k) = \sigma\abs{k} e_\sigma(k)$;
\item Parity under reflection:
      $e_\sigma(-k) = e_\sigma(k)$;
\item Conjugation:
      $\overline{e_\sigma(k)} = e_{-\sigma}(k)$.
\end{enumerate}
\end{proposition}

\begin{proof}
(i) Direct computation using the reality and orthonormality of 
$e_1,e_2$. (ii) Immediate from $e_i\cdot k = 0$. (iii) Using 
$\hat k\times e_1 = e_2$ and $\hat k\times e_2 = -e_1$, we compute
\begin{equation*}
\ii\, k\times e_+(k) 
= \frac{\ii \abs{k}}{\sqrt 2}\,(\hat k\times e_1 + \ii \hat k\times e_2)
= \frac{\ii\abs{k}}{\sqrt 2}\,(e_2 - \ii e_1)
= \abs{k}\,\frac{e_1 + \ii e_2}{\sqrt 2} 
= \abs{k} e_+(k),
\end{equation*}
and similarly for $e_-$. (iv) Immediate from $e_i(-k) = e_i(k)$ and 
$\hat{-k} = -\hat k$ combined with the definition. (v) Direct computation.
\end{proof}

\begin{remark}\label{part-1:rem:helicity-physical}
The vectors $e_\pm(k)$ are eigenvectors of the curl operator with 
eigenvalues $\pm|k|$. They carry the two possible \emph{helicities} of 
a plane wave with wavevector $k$: $e_+$ corresponds to positive 
helicity, $e_-$ to negative helicity. The terminology is that of 
Waleffe \cite{Waleffe1992} and Cambon--Cole.
\end{remark}

\subsection{The state space}
\label{part-1:ssec:state}

We decompose the Fourier coefficient of the velocity as
\begin{equation}\label{part-1:eq:decomposition}
\hat u(k) = u_+(k)\,e_+(k) + u_-(k)\,e_-(k),
\qquad u_\pm(k)\in\C,
\end{equation}
for $k\in\ZZ^3_\ast$. The reality condition \eqref{part-1:eq:reality} becomes
\begin{equation}\label{part-1:eq:reality-helical}
u_\sigma(-k) = \overline{u_\sigma(k)},
\qquad \sigma\in\{+,-\}.
\end{equation}

\begin{notation}[Compound index]\label{part-1:not:compound}
We denote by $\alpha = (k,\sigma)$ a \emph{mode}, with 
$k\in\ZZ^3_\ast$ and $\sigma\in\{+,-\}$, and we write
\begin{equation}\label{part-1:eq:mode-set}
\KK := \ZZ^3_\ast \times \{+,-\}.
\end{equation}
For a finite truncation at wavenumber $K_{\max}$, we write
\begin{equation}\label{part-1:eq:truncated-set}
\KK_{K_{\max}} := \{(k,\sigma)\in\KK : \abs{k}\leq K_{\max}\}.
\end{equation}
The corresponding spectral state is
\begin{equation}\label{part-1:eq:state}
U := (U_\alpha)_{\alpha\in\KK_{K_{\max}}} \in \C^{2\abs{\KK_{K_{\max}}}},
\qquad
U_{(k,\sigma)} := u_\sigma(k).
\end{equation}
We denote by
\begin{equation}\label{part-1:eq:H-s-norm-state}
\norm{U}_{H^s}^2 := \sum_{\alpha\in\KK_{K_{\max}}}
(1 + \abs{k_\alpha}^2)^s \abs{U_\alpha}^2
\end{equation}
the truncated $H^s$ norm.
\end{notation}

The energy of the field $u$ reads, in terms of $U$,
\begin{equation}\label{part-1:eq:energy-state}
E(U) = \frac{1}{2}\norm{u}_{L^2}^2 = \frac{1}{2}\sum_{\alpha} \abs{U_\alpha}^2,
\end{equation}
and the helicity
\begin{equation}\label{part-1:eq:helicity-state}
\HH(U) = \frac{1}{2}\sum_{k\in\ZZ^3_\ast} \abs{k}
\bigl(\abs{u_+(k)}^2 - \abs{u_-(k)}^2\bigr).
\end{equation}

% ============================================================
% SECTION 3 — THE TRIADIC TENSOR
% ============================================================
\section{The triadic tensor}\label{part-1:sec:tensor}

\subsection{Definition}
\label{part-1:ssec:tensor-def}

Substituting the decomposition \eqref{part-1:eq:decomposition} into 
\eqref{part-1:eq:NS-Fourier} and projecting onto the helicoidal basis, we 
obtain the following tensorial representation of the nonlinearity.

\begin{definition}[Triadic tensor]\label{part-1:def:tensor}
For $\gamma = (k,\sigma)$, $\alpha = (p,\sigma')$, 
$\beta = (q,\sigma'')\in\KK$, we define the \emph{triadic tensor}
\begin{equation}\label{part-1:eq:tensor}
\boxed{
\Tt_{\gamma\alpha\beta}
= -2\pi\ii\,
\bigl[e_{\sigma'}(p)\cdot q\bigr]\,
\bigl[\overline{e_\sigma(k)}\cdot e_{\sigma''}(q)\bigr]\,
\delta_{k,p+q}.
}
\end{equation}
\end{definition}

\begin{remark}\label{part-1:rem:tensor-real-form}
The tensor $\Tt_{\gamma\alpha\beta}$ differs from the physical 
Navier--Stokes nonlinearity by the factor $-2\pi\ii$; this convention 
absorbs the imaginary unit into the tensor, so that the algebraic 
properties below take their cleanest form.
\end{remark}

\subsection{Exact factorisation}
\label{part-1:ssec:factorisation}

The tensor \eqref{part-1:eq:tensor} admits an exact factorisation into a 
magnitude and a geometric transition factor.

\begin{definition}[Helical symbols]\label{part-1:def:symbols}
We introduce the \emph{amplitude symbol}
\begin{equation}\label{part-1:eq:amplitude-symbol}
A_{\sigma'}(p,q) := e_{\sigma'}(p)\cdot q,
\end{equation}
and the \emph{transition symbol}
\begin{equation}\label{part-1:eq:transition-symbol}
B_{\sigma\sigma''}(k,q) := \overline{e_\sigma(k)}\cdot e_{\sigma''}(q).
\end{equation}
\end{definition}

\begin{proposition}[Factorisation]\label{part-1:prop:factorisation}
The triadic tensor factorises as
\begin{equation}\label{part-1:eq:factorisation}
\boxed{
\Tt_{\gamma\alpha\beta}
= -2\pi\ii\, A_{\sigma'}(p,q)\, B_{\sigma\sigma''}(k,q)\, 
\delta_{k,p+q}.
}
\end{equation}
\end{proposition}

\begin{proof}
Immediate from \eqref{part-1:eq:tensor} and the definitions.
\end{proof}

\begin{lemma}[Hermitian structure of the transition symbol]
\label{part-1:lem:transition-symbol}
For every $k,q\in\ZZ^3_\ast$,
\begin{equation}\label{part-1:eq:B-hermitian}
B_{\sigma\sigma''}(k,q) = \overline{B_{\sigma''\sigma}(q,k)},
\qquad
\sum_{\sigma''\in\{+,-\}} \abs{B_{\sigma\sigma''}(k,q)}^2 
= 1 - \abs{\hat k\cdot e_\sigma(q)}^2.
\end{equation}
\end{lemma}

\begin{proof}
The first identity follows from the conjugation property 
$\overline{e_\sigma} = e_{-\sigma}$ and the antilinearity of the 
Hermitian inner product. For the second, note that 
$\{e_+(q),e_-(q)\}$ is an orthonormal basis of the plane 
$q^\perp$; the identity is the Parseval formula for the orthogonal 
projection of $e_\sigma(k)$ onto $q^\perp$.
\end{proof}

\subsection{The eight helicoidal channels}
\label{part-1:ssec:channels}

The tensor indices $\sigma,\sigma',\sigma''\in\{+,-\}$ partition 
$\Tt_{\gamma\alpha\beta}$ into \emph{eight helicoidal channels}, 
indexed by the triple $(\sigma,\sigma',\sigma'')$.

\begin{definition}[Channels]\label{part-1:def:channels}
We call the two channels
\begin{equation}\label{part-1:eq:pure-channels}
\Tt^+ := \Tt_{\gamma\alpha\beta}\big|_{\sigma\sigma'\sigma'' = +++},
\qquad
\Tt^- := \Tt_{\gamma\alpha\beta}\big|_{\sigma\sigma'\sigma'' = ---},
\end{equation}
the \emph{pure channels}. The six remaining channels are called 
\emph{acopling channels}, and we group them as
\begin{equation}\label{part-1:eq:ac-channels}
\Tt^{\mathrm{ac}} 
:= \bigl\{\Tt^{\sigma\sigma'\sigma''} : \sigma\sigma'\sigma'' \neq +++ \text{ and } \neq ---\bigr\}.
\end{equation}
\end{definition}

\begin{remark}\label{part-1:rem:ac-two-classes}
Among the six acopling channels, the three with two equal helicities 
and one opposite ($++-$, $+-+$, $-++$) form the \emph{2-against-1} 
class; the other three ($+--$, $-+-$, $--+$) form the 
\emph{1-against-2} class. This distinction will be relevant in the 
sequels.
\end{remark}

\subsection{The spectral evolution equation}
\label{part-1:ssec:spectral-eq}

In terms of the tensor \eqref{part-1:eq:tensor} and the state 
$U\in\C^{2\abs{\KK_{K_{\max}}}}$, the truncated Navier--Stokes equations 
read
\begin{equation}\label{part-1:eq:spectral-evolution}
\boxed{
\partial_t U_\gamma 
= \sum_{\alpha,\beta\in\KK_{K_{\max}}} \Tt_{\gamma\alpha\beta}\, U_\alpha U_\beta
- \nu\abs{k_\gamma}^2 U_\gamma 
+ F_\gamma.
}
\end{equation}

Introducing the notation
\begin{equation}\label{part-1:eq:N-def}
N_\gamma(U) := \sum_{\alpha,\beta} \Tt_{\gamma\alpha\beta} U_\alpha U_\beta,
\qquad
D := \operatorname{diag}(\abs{k_\gamma}^2),
\end{equation}
the evolution \eqref{part-1:eq:spectral-evolution} writes compactly as
\begin{equation}\label{part-1:eq:spectral-compact}
\partial_t U = -N(U) - \nu D U + F.
\end{equation}

\begin{remark}[Sign convention]\label{part-1:rem:sign}
The minus sign in front of $N(U)$ anticipates the structural 
interpretation of $N$ as the bracket of a Lie--Poisson algebra; see 
Section \ref{part-1:ssec:LiePoisson}.
\end{remark}

\subsection{Bivariant Lie Algebra and Geometric Saturation}
\label{part-1:ssec:lean-lie}

The factorisation of the triadic tensor \eqref{part-1:eq:factorisation} and the spectral evolution \eqref{part-1:eq:spectral-compact} imply that the non-linear advection term $N(U)$ is strictly isomorphic to a bivariant Lie commutator $[\Psi, \Psi]$. This elevates the classical Leray projection to a strict algebraic bracket, ensuring that the inviscid limits natively preserve orthogonal conservation.

\begin{theorem}[Symplectic Phase Saturation]\label{part-1:thm:saturation}
The dynamic braking inequality dictates that the heterochiral (stretching) growth rate induces a homochiral braking such that $\gamma_{\text{homo}} \leq -2 |\gamma_{\text{hetero}}|$. This geometric interaction enforces a strict phase desynchronization anomaly tied to the amplitude, known as the \emph{Saturn Symplectic Invariant}. 
\end{theorem}

\begin{corollary}[Zero Trace Condition]\label{part-1:cor:zero-trace}
The mathematical foundation of the Saturn Symplectic Invariant rigorously establishes the exact orthogonality of the non-linear operator, yielding the Zero Trace constraint:
\begin{equation}
[\Psi, \Psi] \equiv 0
\end{equation}
This condition strictly prohibits finite-time blow-up, bounding the coherence functional $\chi(\rho)$ away from $1$ and forcing the system to dissipate its phase energy through ordered geometric channels.
\end{corollary}

This foundational structure, formally verified in the Lean 4 theorem prover, is what mathematically enables the geometric saturation bounds and dimensional collapse explored in the subsequent parts.

% ============================================================
% SECTION 4 — ALGEBRAIC IDENTITIES
% ============================================================
\section{Algebraic identities}\label{part-1:sec:identities}

Throughout this section, the truncation $K_{\max}<\infty$ is fixed, and 
we write $\KK$ for $\KK_{K_{\max}}$.

\subsection{Conjugate-parity relation}
\label{part-1:ssec:parity}

\begin{theorem}[Conjugate-parity]\label{part-1:thm:parity}
For every $\gamma,\alpha,\beta\in\KK$,
\begin{equation}\label{part-1:eq:parity-1}
\boxed{
\Tt_{-\gamma,-\alpha,-\beta} = -\,\Tt_{\gamma\alpha\beta},
}
\end{equation}
where $-\gamma := (-k_\gamma, \sigma_\gamma)$, and similarly for 
$-\alpha,-\beta$. Equivalently,
\begin{equation}\label{part-1:eq:parity-2}
\overline{\Tt_{-\gamma,-\alpha,-\beta}} = \Tt_{\gamma\alpha\beta}.
\end{equation}
\end{theorem}

\begin{proof}
Recall the factorisation \eqref{part-1:eq:factorisation}. Under the 
reflections $k\to -k$, $p\to -p$, $q\to -q$, we have:
\begin{align*}
A_{\sigma'}(-p,-q) &= e_{\sigma'}(-p)\cdot(-q) 
= -\,e_{\sigma'}(p)\cdot q 
= -A_{\sigma'}(p,q), \\
B_{\sigma\sigma''}(-k,-q) &= \overline{e_\sigma(-k)}\cdot e_{\sigma''}(-q)
= \overline{e_\sigma(k)}\cdot e_{\sigma''}(q) 
= B_{\sigma\sigma''}(k,q), \\
\delta_{-k,-p-q} &= \delta_{k,p+q}.
\end{align*}
Therefore
\begin{equation*}
\Tt_{-\gamma,-\alpha,-\beta} 
= -2\pi\ii\, \bigl(-A_{\sigma'}(p,q)\bigr) B_{\sigma\sigma''}(k,q) \delta_{k,p+q}
= -\,\Tt_{\gamma\alpha\beta}.\qedhere
\end{equation*}
\end{proof}

\begin{corollary}[Channel pairings]\label{part-1:cor:channel-pair}
The eight channels pair up under conjugate-parity as
\begin{equation}\label{part-1:eq:channel-pairings}
\Tt^{---} = -\,\overline{\Tt^{+++}},
\qquad
\Tt^{--+} = -\,\overline{\Tt^{++-}},
\qquad
\Tt^{-+-} = -\,\overline{\Tt^{+-+}},
\qquad
\Tt^{-++} = -\,\overline{\Tt^{+--}}.
\end{equation}
Consequently, there are exactly four independent complex channels.
\end{corollary}

\begin{proof}
Apply Theorem \ref{part-1:thm:parity} in conjunction with the 
conjugation property of the symbols; the four identities follow by 
direct computation on each of the four classes.
\end{proof}

\subsection{Energy conservation}
\label{part-1:ssec:energy-cons}

\begin{theorem}[Energy conservation]\label{part-1:thm:energy}
Let $U(t)$ solve \eqref{part-1:eq:spectral-compact} with $\nu = 0$ and $F = 0$, 
and with initial datum satisfying the reality condition 
\eqref{part-1:eq:reality-helical}. Then
\begin{equation}\label{part-1:eq:energy-conservation}
\frac{\dd}{\dd t} E(U(t)) = 0
\qquad\text{for all }t\geq 0.
\end{equation}
\end{theorem}

\begin{proof}
We compute
\begin{equation*}
\frac{\dd}{\dd t} E(U) 
= \mathrm{Re}\sum_\gamma \overline{U_\gamma}\,\dot U_\gamma
= -\mathrm{Re}\sum_{\gamma,\alpha,\beta} 
\overline{U_\gamma}\,\Tt_{\gamma\alpha\beta}\, U_\alpha U_\beta.
\end{equation*}
Denote by $S$ the triple sum appearing above. Under the 
involutive change of variables $(\gamma,\alpha,\beta)\mapsto(-\gamma,-\alpha,-\beta)$, 
using Theorem \ref{part-1:thm:parity} and the reality condition, each 
summand transforms as
\begin{equation*}
\overline{U_{-\gamma}}\,\Tt_{-\gamma,-\alpha,-\beta}\,U_{-\alpha} U_{-\beta}
= -\overline{\overline{U_\gamma}\,\Tt_{\gamma\alpha\beta}\,U_\alpha U_\beta}.
\end{equation*}
Hence $S = -\overline{S}$, so $S\in \ii\R$, and consequently 
$\mathrm{Re}(S) = 0$. The conclusion follows.
\end{proof}

\begin{corollary}[Reality preservation]\label{part-1:cor:reality}
Under the hypotheses of Theorem \ref{part-1:thm:energy}, the reality condition 
$U_{-\gamma}(t) = \overline{U_\gamma(t)}$ is preserved for all $t$.
\end{corollary}

\begin{proof}
Let $V_\gamma := U_{-\gamma} - \overline{U_\gamma}$. The evolution of 
$V$ is linear in $V$ with coefficients depending on $U$, and $V(0) = 0$. 
By uniqueness of the ordinary differential equation, $V\equiv 0$.
\end{proof}

\subsection{Channel decoupling}
\label{part-1:ssec:decoupling}

\begin{theorem}[Channel decoupling]\label{part-1:thm:decoupling}
For every $U$ satisfying the reality condition,
\begin{equation}\label{part-1:eq:decoupling}
\boxed{
\mathrm{Re}\ip{U}{N^{\mathrm{pure}}(U)} = 0,
\qquad
\mathrm{Re}\ip{U}{N^{\mathrm{ac}}(U)} = 0,
}
\end{equation}
where $\ip{U}{V} := \sum_\gamma \overline{U_\gamma} V_\gamma$ and
$N^{\mathrm{pure}}$, $N^{\mathrm{ac}}$ denote the contributions to 
$N(U)$ from the pure and acopling channels respectively.
\end{theorem}

\begin{proof}
The argument of Theorem \ref{part-1:thm:energy} applies channel by channel. 
Conjugate-parity pairs each channel with its opposite; the cancellation 
occurs independently within each pair. Hence each channel-family 
contributes a purely imaginary scalar to the sum, and its real part 
vanishes. 
\end{proof}

\begin{remark}\label{part-1:rem:decoupling-meaning}
Theorem \ref{part-1:thm:decoupling} states that the pure and acopling 
contributions to the energy balance vanish \emph{separately}, not merely 
in sum. This is a strictly stronger statement than the standard energy 
conservation and reflects the helical structure of the equations.
\end{remark}

\subsection{Hermitian orthogonality}
\label{part-1:ssec:hermitian}

\begin{definition}[Linearised operator]\label{part-1:def:L}
For fixed $U\in\C^{2\abs{\KK}}$, the \emph{linearised operator} 
$\LL[U]$ is defined by
\begin{equation}\label{part-1:eq:L-def}
(\LL[U] V)_\gamma 
:= \sum_{\alpha,\beta} \bigl(\Tt_{\gamma\alpha\beta} + \Tt_{\gamma\beta\alpha}\bigr)\, 
U_\beta V_\alpha,
\qquad V\in\C^{2\abs{\KK}}.
\end{equation}
\end{definition}

The operator $\LL[U]$ is the Fr\'echet derivative of $N$ at $U$, 
i.e., $N(U+\eps V) = N(U) + \eps \LL[U] V + O(\eps^2)$.

\begin{definition}[Hermitian decomposition]\label{part-1:def:H-A}
We define
\begin{equation}\label{part-1:eq:HA}
H(U) := \tfrac{1}{2}\bigl(\LL[U] + \LL[U]^\dagger\bigr),
\qquad
A(U) := \tfrac{1}{2}\bigl(\LL[U] - \LL[U]^\dagger\bigr),
\end{equation}
so that $\LL[U] = H(U) + A(U)$ with $H$ hermitian and $A$ anti-hermitian.
\end{definition}

\begin{theorem}[Hermitian orthogonality]\label{part-1:thm:H-orth}
For every $U$ satisfying the reality condition,
\begin{equation}\label{part-1:eq:H-orth}
\boxed{
U^\dagger H(U) U = 0.
}
\end{equation}
Equivalently, $\mathrm{Re}\bigl(U^\dagger \LL[U] U\bigr) = 0$.
\end{theorem}

\begin{proof}
By Definition \ref{part-1:def:H-A} and the hermiticity of $H$,
\begin{equation*}
U^\dagger H(U) U 
= \mathrm{Re}\bigl(U^\dagger \LL[U] U\bigr).
\end{equation*}
Expanding,
\begin{equation*}
U^\dagger \LL[U] U 
= \sum_{\gamma,\alpha,\beta} \overline{U_\gamma}
\bigl(\Tt_{\gamma\alpha\beta} + \Tt_{\gamma\beta\alpha}\bigr) U_\beta U_\alpha
= 2S,
\end{equation*}
where $S$ is the triple sum of the proof of Theorem \ref{part-1:thm:energy}. 
The reality argument gives $S\in\ii\R$, hence 
$U^\dagger H(U) U = \mathrm{Re}(2S) = 0$.
\end{proof}

\begin{corollary}[Evolution balance]\label{part-1:cor:balance}
Along a solution of \eqref{part-1:eq:spectral-compact},
\begin{equation}\label{part-1:eq:balance}
\frac{\dd}{\dd t}\norm{U}^2 
= -2\nu\ip{U}{DU} + 2\,\mathrm{Re}\ip{U}{F}.
\end{equation}
In particular, the nonlinearity does not contribute to the evolution 
of the $L^2$ norm.
\end{corollary}

\begin{proof}
Write $\dot U = -H(U)U - A(U)U - \nu DU + F$. Multiplying by 
$U^\dagger$ and taking real parts, the hermitian term vanishes by 
Theorem \ref{part-1:thm:H-orth}, and the anti-hermitian term contributes a 
purely imaginary scalar. The conclusion follows.
\end{proof}

% ============================================================
% SECTION 5 — THE PER-TRIAD RATIO
% ============================================================
\section{The per-triad ratio}\label{part-1:sec:triad}

\subsection{Explicit form of the symbols}
\label{part-1:ssec:symbols-explicit}

We compute the transition symbol $B_{\sigma\sigma''}(k,q)$ explicitly 
in terms of the angle $\theta := \angle(k,q)$.

\begin{lemma}[Explicit transition symbol]\label{part-1:lem:B-explicit}
Let $k,q\in\ZZ^3_\ast$ with $\abs{k},\abs{q}\neq 0$, and let 
$\theta$ be the angle between them. Then, in the gauge 
$e_1(k) = z\text{-axis aligned}$, one has
\begin{equation}\label{part-1:eq:B-explicit}
\abs{B_{++}(k,q)}^2 = \frac{1}{4}(1+\cos\theta)^2,
\qquad
\abs{B_{+-}(k,q)}^2 = \frac{1}{4}(1-\cos\theta)^2.
\end{equation}
Consequently,
\begin{equation}\label{part-1:eq:B-sum}
\abs{B_{++}(k,q)}^2 + \abs{B_{+-}(k,q)}^2 
= \frac{1}{2}(1 + \cos^2\theta).
\end{equation}
\end{lemma}

\begin{proof}
Choose an orthonormal frame in which 
$\hat k = (0,0,1)$ and 
$\hat q = (\sin\theta, 0, \cos\theta)$. Then
\begin{equation*}
e_+(k) = \frac{(1,\ii,0)}{\sqrt 2},
\qquad
e_\pm(q) = \frac{(\cos\theta, \pm\ii, -\sin\theta)}{\sqrt 2}.
\end{equation*}
Direct computation of the Hermitian products gives
\begin{equation*}
\overline{e_+(k)}\cdot e_+(q) = \frac{1}{2}(1+\cos\theta),
\qquad
\overline{e_+(k)}\cdot e_-(q) = \frac{1}{2}(\cos\theta-1).
\end{equation*}
Taking squared moduli yields the stated identities. The sum follows 
immediately.
\end{proof}

\begin{remark}\label{part-1:rem:complete}
The sum \eqref{part-1:eq:B-sum} is exactly 
$1 - \abs{\hat k\cdot e_+(q)}^2$, as predicted by the general 
completeness identity of Lemma \ref{part-1:lem:transition-symbol}. In the 
present gauge, $\hat k\cdot e_+(q) = \sin\theta/\sqrt 2$, so 
$\abs{\hat k\cdot e_+(q)}^2 = \sin^2\theta/2$, and 
$1 - \sin^2\theta/2 = (1+\cos^2\theta)/2$, in agreement.
\end{remark}

\subsection{The per-triad ratio}
\label{part-1:ssec:ratio}

\begin{definition}[Per-triad ratio]\label{part-1:def:ratio}
For a triadic configuration $(k,p,q)$ with $k = p+q$, define the 
\emph{per-triad ratio}
\begin{equation}\label{part-1:eq:ratio-def}
R_{kpq}^2 
:= \frac{\abs{\Tt^{++-}_{kpq}}^2 + \abs{\Tt^{+-+}_{kpq}}^2}
        {\abs{\Tt^{+++}_{kpq}}^2 + \abs{\Tt^{---}_{kpq}}^2}
= \frac{\abs{\Tt^{\mathrm{ac}}_{kpq}}^2}{\abs{\Tt^{\mathrm{pure}}_{kpq}}^2}.
\end{equation}
\end{definition}

\begin{theorem}[Exact form of the ratio]\label{part-1:thm:ratio-exact}
For every triadic configuration $(k,p,q)$ with $k=p+q$, the 
per-triad ratio admits the closed expression
\begin{equation}\label{part-1:eq:ratio-exact}
\boxed{
R_{kpq}^2 = 1 + 2\left(\frac{1-\cos\theta}{1+\cos\theta}\right)^2,
}
\end{equation}
where $\theta := \angle(k,q)$.
\end{theorem}

\begin{proof}
We work in the gauge of Lemma \ref{part-1:lem:B-explicit}. The amplitude 
symbol $A_{\sigma'}(p,q) = e_{\sigma'}(p)\cdot q$ does not depend on 
$(\sigma_g,\sigma_b)$; it factors out of the ratio. What remains is 
the ratio of transition symbols:
\begin{equation*}
\frac{\abs{B_{+-}(k,q)}^2 + \abs{B_{-+}(k,q)}^2}
     {\abs{B_{++}(k,q)}^2 + \abs{B_{--}(k,q)}^2}
= \frac{2\abs{B_{+-}(k,q)}^2}{2\abs{B_{++}(k,q)}^2}
= \left(\frac{1-\cos\theta}{1+\cos\theta}\right)^2.
\end{equation*}
Adding the direct channel contribution ($R^2 = 1$ from the pure 
channels alone) yields
\begin{equation*}
R_{kpq}^2 = 1 + 2\left(\frac{1-\cos\theta}{1+\cos\theta}\right)^2,
\end{equation*}
as claimed.
\end{proof}

\begin{corollary}[Lower bound]\label{part-1:cor:ratio-lower}
For every triadic configuration $(k,p,q)$,
\begin{equation}\label{part-1:eq:ratio-lower}
\boxed{
R_{kpq} \geq 1.
}
\end{equation}
Equality holds if and only if $\cos\theta = 1$, i.e., if and only if 
$q\parallel k$.
\end{corollary}

\begin{proof}
The right-hand side of \eqref{part-1:eq:ratio-exact} is $\geq 1$ because the 
squared bracket is non-negative. Equality occurs precisely when 
$1-\cos\theta = 0$.
\end{proof}

\begin{corollary}[No upper bound]\label{part-1:cor:ratio-upper}
\begin{equation}\label{part-1:eq:ratio-upper}
\sup_{k,p,q} R_{kpq} = \infty.
\end{equation}
The supremum is attained in the limit of retrogradely aligned 
configurations, $q\to -k$.
\end{corollary}

\begin{proof}
As $\cos\theta \to -1$, the bracket in \eqref{part-1:eq:ratio-exact} diverges 
as $(1-\cos\theta)^2/(1+\cos\theta)^2 \sim 4/(1+\cos\theta)^2$. Hence 
$R_{kpq}\to\infty$.
\end{proof}

\begin{remark}\label{part-1:rem:ratio-meaning}
Corollary \ref{part-1:cor:ratio-lower} is the exact statement that, for every 
triadic configuration compatible with the Navier--Stokes equations, the 
acopling channels carry at least as much squared amplitude as the pure 
channels. Corollary \ref{part-1:cor:ratio-upper} shows that this dominance is 
unbounded in the retrograde limit; the energy in such configurations 
is however cut off by the Leray projection, which vanishes when 
$q\parallel k$.
\end{remark}

\subsection{The channel hierarchy}
\label{part-1:ssec:hierarchy}

Lemma \ref{part-1:lem:B-explicit} also reveals a hierarchy within the 
acopling channels.

\begin{proposition}[Hierarchy]\label{part-1:prop:hierarchy}
The acopling channels split into two classes with distinct 
$L^2$-norms:
\begin{enumerate}[label=(\roman*),leftmargin=*]
\item The \emph{2-against-1} class, containing $\Tt^{++-}$ and its 
      conjugate $\Tt^{--+}$, has squared norm
      \begin{equation}\label{part-1:eq:hierarchy-1}
      \norm{\Tt^{++-}}_{L^2}^2 \propto \frac{1}{4}(1-\cos\theta)^2;
      \end{equation}
\item The \emph{1-against-2} class, containing $\Tt^{+--}$ and its 
      conjugate $\Tt^{-++}$, has a strictly smaller norm, with ratio 
      \begin{equation}\label{part-1:eq:hierarchy-2}
      \frac{\norm{\Tt^{++-}}}{\norm{\Tt^{+--}}} \approx 6.13
      \end{equation}
      in the truncation $K_{\max} = 5$ with Kolmogorov spectral weight.
\end{enumerate}
\end{proposition}

\begin{proof}
The first statement follows from the explicit computation of 
$B_{\sigma\sigma''}$; the second is obtained by numerical evaluation 
of the $L^2$ norms in the truncation and is reported here for 
reference (see Section \ref{part-1:ssec:numeric-hierarchy}).
\end{proof}

\begin{remark}\label{part-1:rem:hierarchy-physical}
Physically, the dominance of the 2-against-1 channels means that 
interactions in which two waves of one helicity conspire to produce 
one wave of the opposite helicity are the strongest nonlinear 
processes at the level of individual triads. This is a structural 
feature of the Euler/NS tensor in the helical basis and will play a 
central role in the sequel on the renormalisation flow.
\end{remark}

\subsection{The numerical verification}
\label{part-1:ssec:numeric-hierarchy}

The identities of Theorems \ref{part-1:thm:parity}--\ref{part-1:thm:ratio-exact} 
have been verified numerically in the truncations $K_{\max}\in\{2,3,4,5\}$, 
with relative errors below $10^{-15}$ for the algebraic identities and 
below $10^{-13}$ for the per-triad ratio. The channel norms are 
reported in Table \ref{part-1:tab:channel-norms}.

\begin{table}[h]
\centering
\caption{Channel norms in $K_{\max}=5$ with unit energy. The 
$L^2$-norm $\norm{\Tt^\sigma}$ is computed as the Frobenius norm of 
the corresponding tensor block.}
\label{part-1:tab:channel-norms}
\begin{tabular}{lcc}
\toprule
Channel & Structure & $\norm{\cdot}_{L^2}$ \\
\midrule
$+++$ & Pure & $28.19$ \\
$---$ & Pure & $28.19$ \\
$++-$ & Acopling (2-1) & $50.97$ \\
$+-+$ & Acopling (2-1) & $50.97$ \\
$-+-$ & Acopling (2-1) & $50.97$ \\
$--+$ & Acopling (2-1) & $50.97$ \\
$+--$ & Acopling (1-2) & $8.31$ \\
$-++$ & Acopling (1-2) & $8.31$ \\
\bottomrule
\end{tabular}
\end{table}

\begin{remark}\label{part-1:rem:frac-ac}
Table \ref{part-1:tab:channel-norms} shows that 
$\norm{\Tt^{\mathrm{ac}}}_{L^2}/\norm{\Tt}_{L^2} \approx 0.932$, i.e., 
$86.9\%$ of the total Frobenius norm is concentrated in the acopling 
channels. The pure channels carry only $13.1\%$.
\end{remark}

% ============================================================
% SECTION 6 — LIE-POISSON STRUCTURE
% ============================================================
\section{Lie--Poisson structure}
\label{part-1:sec:structure}

\subsection{The Lie algebra}
\label{part-1:ssec:lie-algebra}

\begin{definition}[Bracket]\label{part-1:def:bracket}
For $v_a, v_b \in H^s_{\mathrm{df}}(\T^3)$, define
\begin{equation}\label{part-1:eq:bracket}
[v_a, v_b] := -\Leray\bigl[(v_a\cdot\nabla)v_b - (v_b\cdot\nabla)v_a\bigr].
\end{equation}
\end{definition}

\begin{proposition}\label{part-1:prop:lie-algebra}
The bracket \eqref{part-1:eq:bracket} endows $H^s_{\mathrm{df}}(\T^3)$ with 
the structure of a Lie algebra: it is bilinear, antisymmetric, and 
satisfies the Jacobi identity
\begin{equation}\label{part-1:eq:Jacobi}
[v_a, [v_b, v_c]] + [v_b, [v_c, v_a]] + [v_c, [v_a, v_b]] = 0.
\end{equation}
\end{proposition}

\begin{proof}
Bilinearity and antisymmetry are immediate. The Jacobi identity follows 
from the observation that $[v_a, v_b]$ is the Lie bracket of the 
vector fields $v_a, v_b$, projected onto the divergence-free fields, 
and from the fact that $\Leray$ is a Lie-algebra homomorphism on the 
subalgebra of divergence-free fields.
\end{proof}

In the Fourier representation, the structure constants of this Lie 
algebra coincide, up to the factor $-2\pi\ii$, with the triadic tensor 
$\Tt_{\gamma\alpha\beta}$:
\begin{equation}\label{part-1:eq:structure-constants}
[\hat v_a, \hat v_b]_\gamma = \sum_{\alpha,\beta} \Tt_{\gamma\alpha\beta} 
\,(\hat v_a)_\alpha (\hat v_b)_\beta.
\end{equation}

\subsection{Poisson structure and Casimirs}
\label{part-1:ssec:LiePoisson}

The Navier--Stokes nonlinearity is the coadjoint action of this Lie 
algebra.

\begin{proposition}[Lie--Poisson formulation]\label{part-1:prop:LiePoisson}
The Euler equations ($\nu = f = 0$) are equivalent to the Lie--Poisson 
system
\begin{equation}\label{part-1:eq:LiePoisson}
\dot U = \Pi(U)\, \nabla H(U),
\end{equation}
where $\Pi(U) := \mathrm{ad}^\ast_U$ is the coadjoint operator and 
$H(U) := \tfrac12\norm{U}_{L^2}^2$ is the energy.
\end{proposition}

\begin{proof}
We compute
\begin{equation*}
\Pi(U)\nabla H(U)_\gamma = \sum_{\alpha,\beta} \Tt_{\gamma\alpha\beta} 
U_\alpha U_\beta = N_\gamma(U),
\end{equation*}
which coincides with $-\dot U_\gamma$ by \eqref{part-1:eq:spectral-compact} 
with $\nu = f = 0$.
\end{proof}

\begin{proposition}[Casimirs]\label{part-1:prop:Casimirs}
The functionals
\begin{equation}\label{part-1:eq:Casimirs}
C_1(U) := \tfrac12\norm{U}_{L^2}^2,
\qquad
C_2(U) := \tfrac12\sum_{k} \abs{k}\bigl(\abs{u_+(k)}^2 - \abs{u_-(k)}^2\bigr),
\end{equation}
are Casimirs of the Lie--Poisson structure: they Poisson-commute with 
every functional, in particular $\{C_i, H\} = 0$ for $i = 1, 2$.
\end{proposition}

\begin{proof}
For $C_1$, the result is Theorem \ref{part-1:thm:energy}. For $C_2$, the 
Casimir property is a classical result of Arnold--Khesin; the 
algebraic verification in the truncated setting reduces to the 
symmetry of the tensor under the exchange of helicities combined with 
the conservation of helicity, which we omit here for brevity and 
defer to the standard references.
\end{proof}

\begin{remark}\label{part-1:rem:Casimirs-meaning}
$C_1$ is the energy, $C_2$ is the helicity. The existence of these two 
Casimirs restricts the dynamics of the Euler flow to a codimension-$2$ 
submanifold of the state space, foliated by coadjoint orbits.
\end{remark}

% ============================================================
% APPENDIX
% ============================================================
